% =====================================================================
\chapter{Analyse fichier par fichier}
\label{ch:fichiers}

Chaque fiche suit la même grille : objectif, fonctions, classes, dépendances,
entrées, sorties, fichiers appelés, fichiers appelants, explication du code, et
ce qu'il faut en retenir pour la présentation.

% =====================================================================
\section{\file{backend/app/main.py}}

\textbf{Objectif} --- Assembler l'application : middlewares, routeurs, tâches de
fond.

\textbf{Fonctions principales} --- \fn{lifespan} (gestionnaire de contexte
asynchrone), \fn{\_prime\_storage\_health}, \fn{root}.

\textbf{Classes principales} --- Aucune ; l'objet \texttt{app = FastAPI(...)} est
créé au niveau du module.

\textbf{Dépendances} --- FastAPI, \cls{CORSMiddleware}, \cls{AuditMiddleware},
les 11 routeurs, \file{dashboard/kpi\_service.py},
\file{services/incident\_watcher.py}, \file{services/metrics\_collector.py},
\file{app/models} (import global qui enregistre tous les modèles).

\textbf{Entrées} --- Requêtes HTTP. \textbf{Sorties} --- Réponses HTTP ; la
racine \route{GET /} renvoie \texttt{\{"status": "OASIS AI Copilot backend
running"\}}.

\textbf{Fichiers appelés} --- Tous les routeurs, les deux collecteurs.
\textbf{Fichiers qui l'appellent} --- Uvicorn
(\texttt{uvicorn app.main:app -{}-reload -{}-port 8000}, cf.
\file{.claude/launch.json}).

\textbf{Explication} --- Le \fn{lifespan} crée trois tâches
(\fn{metrics\_collector.run\_forever}, \fn{incident\_watcher.run\_forever},
\fn{\_prime\_storage\_health}), \texttt{yield}, puis les annule et attend leur
\cls{CancelledError} à l'arrêt. \fn{\_prime\_storage\_health} sort le premier
relevé SSH du chemin d'une requête : « sans lui, le premier affichage du
dashboard après un redémarrage attendrait la poignée de main SSH ».

\textbf{À retenir} --- Un backend qui \emph{travaille sans qu'on le lui
demande}. C'est ce qui distingue OASIS d'un simple chatbot : il surveille,
détecte et mesure en continu.

% =====================================================================
\section{\file{backend/app/core/config.py}}

\textbf{Objectif} --- Centraliser la configuration, validée au démarrage.

\textbf{Classes} --- \cls{Settings(BaseSettings)}, instanciée une fois en
\fn{settings}.

\textbf{Dépendances} --- \texttt{pydantic-settings}.
\textbf{Entrées} --- Variables d'environnement et \file{.env}.
\textbf{Sorties} --- L'objet \fn{settings}.

\textbf{Fichiers qui l'appellent} --- \file{security.py}, \file{session.py},
\file{ssh.py}, \file{groq\_client.py}, \file{vectorstore.py},
\file{embedding\_service.py}, \file{knowledge\_service.py},
\file{request\_service.py}.

\textbf{À retenir} --- Les variables SSH sont \emph{optionnelles} par choix, avec
une erreur explicite au moment de l'appel. Le nom du modèle d'embedding est
configurable, et le changer invalide les collections Chroma --- ce que le code
détecte au lieu de laisser passer.

% =====================================================================
\section{\file{backend/app/api/deps.py}}

\textbf{Objectif} --- Fournir les dépendances d'authentification et
d'autorisation.

\textbf{Fonctions} --- \fn{get\_current\_user(credentials, db)},
\fn{require\_role(*allowed\_roles)}.

\textbf{Entrées} --- En-tête \texttt{Authorization: Bearer}, session DB.
\textbf{Sorties} --- Une entité \cls{User}, ou \cls{HTTPException} 401/403.

\textbf{Fichiers appelés} --- \file{database/session.py},
\file{core/security.py}, \file{models/user.py}.
\textbf{Fichiers qui l'appellent} --- Les 11 routeurs.

\textbf{À retenir} --- Deux contrôles, pas un : le jeton \emph{et} l'état du
compte en base. Un compte désactivé est rejeté immédiatement, sans attendre
l'expiration de son jeton.

% =====================================================================
\section{\file{backend/app/api/v1/storage.py}}

\textbf{Objectif} --- Exposer le chat : gestion des fils et tour de dialogue.

\textbf{Fonctions} --- \fn{\_require\_conversation}, \fn{list\_conversations},
\fn{conversation\_messages}, \fn{rename\_conversation},
\fn{delete\_conversation}, \fn{chat} (\texttt{async}).

\textbf{Dépendances} --- \file{services/chat\_service.py},
\file{personas/storage/agent.py}, \file{services/incident\_investigation.py},
\file{services/incident\_service.py}, \file{schemas/chat.py}.

\textbf{Entrées} --- \cls{ChatRequest}, \cls{ConversationUpdate}, identifiants de
fil. \textbf{Sorties} --- \cls{ConversationOut}, \cls{ChatMessageOut},
\cls{ChatTurnOut}, 204.

\textbf{Explication} --- L'ordre des opérations dans \fn{chat} est le point clé :
l'agent tourne \emph{avant} toute écriture. Un échec de Groq ou de la VM
n'ajoute donc ni fil vide ni question sans réponse.

\textbf{À retenir} --- Le 404 sur un fil appartenant à un collègue est un choix de
sécurité : « dire "interdit" plutôt qu'"introuvable" confirmerait son
existence ».

% =====================================================================
\section{\file{backend/app/personas/storage/agent.py}}

\textbf{Objectif} --- Le cerveau : RAG, dialogue avec le LLM, appels d'outils,
garde-fou de sécurité, détection d'anomalie.

\textbf{Classes} --- \cls{IncidentContext}, \cls{AgentTurn},
\cls{StoragePersona} (instanciée en \fn{storage\_persona}).

\textbf{Fonctions} --- \fn{\_candidate\_calls}, \fn{\_recover\_written\_call},
\fn{StoragePersona.\_propose\_action}, \fn{StoragePersona.handle\_message}.

\textbf{Dépendances} --- \file{llm/groq\_client.py}, \file{mcp/storage/client.py},
\file{mcp/storage/server.py} (\texttt{REMEDIATION\_TOOLS}),
\file{personas/storage/incident\_rules.py},
\file{personas/storage/prompts.py}, \file{rag/retriever.py},
\file{services/incident\_service.py}, \file{database/session.py}.

\textbf{Entrées} --- \texttt{user\_message}, \texttt{created\_by},
\texttt{history}, \texttt{incident}.
\textbf{Sorties} --- \cls{AgentTurn}. Effets de bord : lignes
\tbl{mcp\_calls}, \tbl{action\_requests}, \tbl{incident\_tickets}.

\textbf{Fichiers qui l'appellent} --- \file{api/v1/storage.py},
\file{services/incident\_investigation.py}.

\textbf{Constantes} --- \texttt{\_MAX\_TOOL\_ROUNDS = 6},
\texttt{\_TARGET\_ARGUMENTS = ("target", "volume\_id", "job\_id")},
\texttt{\_JSON\_OBJECT} (expression régulière).

\textbf{À retenir} --- Trois idées à exposer : \emph{(1)} le garde-fou qui
transforme un appel d'outil en proposition ; \emph{(2)} la correction
automatique de la cible (libellé \arrowto{} identifiant technique) ; \emph{(3)}
le rattrapage d'un appel écrit en prose, avec double validation.

% =====================================================================
\section{\file{backend/app/personas/storage/prompts.py}}

\textbf{Objectif} --- Le prompt système du persona.

\textbf{Contenu} --- Une constante, \fn{STORAGE\_SYSTEM\_PROMPT}. Elle définit
quatre domaines d'expertise, quatre règles de sûreté (ne jamais prétendre avoir
agi sans résultat d'outil, demander confirmation pour une action irréversible,
séparer faits / hypothèses / recommandations, préférer la documentation interne
et ne jamais citer une documentation absente du contexte), et impose une
réponse en anglais.

\textbf{À retenir} --- Une nuance de rédaction remarquable : la spécialité
« façonne le ton et les exemples ; elle ne limite pas les questions acceptées ».
Le prompt interdit explicitement de refuser une question hors domaine --- un
défaut classique des assistants spécialisés.

% =====================================================================
\section{\file{backend/app/personas/storage/incident\_rules.py}}

\textbf{Objectif} --- Décider, de façon déterministe, ce qui constitue un
incident.

\textbf{Fonctions} --- \fn{\_extract\_items}, \fn{check\_capacity},
\fn{check\_backup\_status}, \fn{evaluate\_tool\_result}.

\textbf{Constantes} --- \texttt{CAPACITY\_WARNING\_THRESHOLD = 85},
\texttt{CAPACITY\_CRITICAL\_THRESHOLD = 95},
\texttt{\_KNOWN\_LIST\_WRAPPER\_KEYS}, \texttt{\_RULES\_BY\_TOOL}.

\textbf{Entrées} --- Nom d'outil et résultat brut.
\textbf{Sorties} --- Liste de dictionnaires prêts pour
\fn{create\_ticket\_from\_diagnosis}.

\textbf{Fichiers qui l'appellent} --- \file{agent.py},
\file{incident\_watcher.py}, \file{incident\_verification.py}.

\textbf{Explication} --- \fn{\_extract\_items} normalise trois formes
d'enveloppe (liste brute, dictionnaire enveloppant une liste connue,
dictionnaire plat). Le docstring documente un bug corrigé :
« \fn{get\_capacity()} retourne \texttt{\{"volumes": [...]\}} \ldots{}
L'ancienne version ne gérait que les deux derniers cas, donc ne détectait jamais
rien sur \fn{get\_capacity}. »

\textbf{À retenir} --- Le même fichier sert à \emph{ouvrir} et à \emph{fermer}
un ticket. C'est ce qui garantit la cohérence : ce qui a déclenché l'alerte est
exactement ce qui doit avoir disparu.

% =====================================================================
\section{\file{backend/app/mcp/storage/server.py}}

\textbf{Objectif} --- Registre des outils et classification du risque.

\textbf{Contenu} --- \fn{TOOL\_REGISTRY} (10 entrées),
\texttt{DESTRUCTIVE\_TOOLS} (1 entrée), \texttt{REMEDIATION\_TOOLS} (4 entrées),
\fn{risk\_for}, \fn{invoke}.

\textbf{Fichiers qui l'appellent} --- \file{tool\_server.py}
(\fn{TOOL\_REGISTRY}), \file{agent.py} et \file{incident\_investigation.py}
(\texttt{REMEDIATION\_TOOLS}).

\textbf{À retenir} --- 60 lignes qui portent toute la politique de sécurité
d'exécution. Signaler honnêtement que \fn{risk\_for} et \fn{invoke} ne sont pas
encore utilisés.

% =====================================================================
\section{\file{backend/app/mcp/storage/tool\_server.py}}

\textbf{Objectif} --- Le serveur MCP lui-même.

\textbf{Contenu} --- 12 lignes : import de \cls{MCPServer}, instanciation,
boucle d'enregistrement, \texttt{mcp.run(transport="stdio")}.

\textbf{Entrées} --- JSON-RPC sur \texttt{stdin}.
\textbf{Sorties} --- JSON-RPC sur \texttt{stdout}.

\textbf{Fichiers appelés} --- \file{server.py} et, transitivement, les sept
modules d'outils. \textbf{Lancé par} --- \file{client.py}, comme sous-processus.

\textbf{À retenir} --- La brièveté est l'argument : le SDK MCP dérive les
schémas des annotations de type et des docstrings, il n'y a donc rien à écrire
à la main.

% =====================================================================
\section{\file{backend/app/mcp/storage/client.py}}

\textbf{Objectif} --- Piloter le serveur MCP et auditer chaque appel.

\textbf{Classes} --- \cls{StorageMCPClient} (protocole de gestionnaire de
contexte asynchrone).

\textbf{Fonctions} --- \fn{start}, \fn{stop}, \fn{\_require\_session},
\fn{list\_tools\_for\_llm}, \fn{call\_tool}.

\textbf{Dépendances} --- \texttt{mcp.ClientSession},
\texttt{mcp.client.stdio}, \file{models/mcp\_call.py},
\file{services/mcp\_audit.py}.

\textbf{Entrées de \fn{call\_tool}} --- \texttt{tool\_name}, puis les mots-clés
\texttt{db}, \texttt{user\_id}, \texttt{correlation\_id},
\texttt{action\_request\_id}, et \texttt{**arguments}.
\textbf{Sorties} --- \texttt{dict} décodé, ou \cls{RuntimeError}.

\textbf{Fichiers qui l'appellent} --- \file{agent.py}, \file{actions.py} ;
\file{incident\_verification.py} reçoit le client déjà démarré.

\textbf{À retenir} --- Pas d'instance partagée au niveau du module : c'est la
correction d'un bug réel, documenté dans \file{actions.py}. Le verrou
\texttt{asyncio.Lock} sérialise les appels, le transport stdio n'étant pas
multiplexé.

% =====================================================================
\section{\file{backend/app/mcp/storage/ssh.py}}

\textbf{Objectif} --- Point d'accès unique à la VM.

\textbf{Classes} --- \cls{StorageVMError(RuntimeError)}.
\textbf{Fonctions} --- \fn{\_client}, \fn{run}, \fn{run\_many}, \fn{quote}.
\textbf{Constantes} --- \texttt{CONNECT\_TIMEOUT\_SECONDS = 5},
\texttt{COMMAND\_TIMEOUT\_SECONDS = 15}.

\textbf{Entrées} --- Une commande, ou un dictionnaire de commandes.
\textbf{Sorties} --- \texttt{stdout} en \texttt{str}, ou \cls{StorageVMError}.

\textbf{Fichiers qui l'appellent} --- Les sept modules d'outils.

\textbf{Explication} --- \fn{\_client} vérifie d'abord que les trois variables
SSH sont renseignées et nomme celles qui manquent. Les messages d'erreur sont
volontairement \emph{en anglais} : « Ces messages sont affichés tels quels dans
l'interface, qui est en anglais. » La connexion est fermée dans un
\texttt{finally}.

\textbf{À retenir} --- 130 lignes qui isolent tout le risque d'exécution. Aucun
outil ne construit sa propre connexion, aucun n'échappe à \fn{quote}.

% =====================================================================
\section{\file{backend/app/mcp/storage/tools/backup.py}}

\textbf{Objectif} --- Lire et piloter les jobs de sauvegarde réels.

\textbf{Fonctions} --- \fn{\_looks\_like\_backup}, \fn{\_parse\_timer\_units},
\fn{\_service\_for}, \fn{\_is\_storage\_backup}, \fn{\_show\_units},
\fn{\_parse\_show}, \fn{\_moment}, \fn{\_hours\_since},
\fn{\_job\_from\_properties}, \fn{list\_backups}, \fn{get\_backup\_job\_status},
\fn{run\_backup}, \fn{get\_backup\_logs}.

\textbf{Sorties} --- \texttt{\{"jobs": [...], "configured": bool\}} ou
\texttt{\{"error": ..., "jobs": []\}}.

\textbf{Trois finesses à connaître}
\begin{enumerate}
  \item \textbf{Filtrage des units système.} \fn{\_is\_storage\_backup} écarte
  \texttt{dpkg-db-backup.timer}, « qui sauvegarde la base de paquets d'Ubuntu »
  : une unit livrée par un paquet n'est retenue que si son nom cite un vrai
  outil de sauvegarde (restic, borg, veeam\ldots).
  \item \textbf{Lecture groupée.} \fn{\_show\_units} concatène toutes les
  commandes en une seule, séparées par le marqueur
  \texttt{\#\#\#OASIS-UNIT=}. Motif : « sur six timers, \fn{list\_backups}
  mettait près de quatre secondes, et le dashboard se figeait d'autant à chaque
  navigation. »
  \item \textbf{« Jamais lancé » n'est pas « a réussi ».} systemd renvoie
  \texttt{Result=success} sur une unit qui n'a jamais démarré. Le code teste
  \texttt{ExecMainStartTimestamp} : sans cela, « un job planifié mais jamais
  exécuté s'afficherait comme une sauvegarde réussie ».
\end{enumerate}

\textbf{À retenir} --- \texttt{configured: False} est distinct de
\texttt{error} : rien n'est configuré, ce n'est pas une panne. Cette distinction
remonte jusqu'au dashboard et jusqu'aux insights.

% =====================================================================
\section{\file{backend/app/mcp/storage/tools/repair.py}}

\textbf{Objectif} --- Réparer un job dont le dépôt restic est introuvable.

\textbf{Fonctions} --- \fn{\_service\_for}, \fn{\_parse\_environment},
\fn{\_parse\_environment\_files}, \fn{\_repository\_state},
\fn{\_healthy\_donor}, \fn{\_drop\_in\_body}, \fn{\_drop\_in\_path},
\fn{\_write\_commands}, \fn{\_write\_drop\_in}, \fn{\_missing\_privilege},
\fn{repair\_backup\_repository}.

\textbf{Entrées} --- \texttt{target}, \texttt{confirm=False}.
\textbf{Sorties} --- Un dictionnaire dont le champ \texttt{action} prend l'une de
sept valeurs : \texttt{repair\_failed}, \texttt{no\_repair\_possible},
\texttt{no\_repair\_needed}, \texttt{repair\_requires\_confirmation},
\texttt{repository\_repaired}.

\textbf{À retenir} --- Le seul outil qui \emph{modifie une configuration}, et
celui qui refuse le plus souvent d'agir. Trois idées à présenter : le
diagnostic du dépôt par un faux mot de passe, le refus en cas d'égalité entre
voisins, et la réversibilité par \texttt{systemctl revert}.

% =====================================================================
\section{\file{backend/app/api/v1/actions.py}}

\textbf{Objectif} --- Le seul endroit du système qui exécute une action sur la
VM.

\textbf{Fonctions} --- \fn{list\_my\_actions}, \fn{propose\_action},
\fn{confirm\_action} (\texttt{async}), \fn{\_close\_the\_loop}
(\texttt{async}), \fn{reject\_action}, \fn{get\_action}.

\textbf{Dépendances} --- \file{mcp/storage/client.py},
\file{services/incident\_service.py}, \file{services/incident\_verification.py},
\file{models/action\_request.py}, \file{schemas/incident.py}.

\textbf{Fichiers qui l'appellent} --- \file{frontend/src/components/ActionProposal.tsx}.

\textbf{À retenir} --- Le docstring de \fn{confirm\_action} raconte un bug de
sécurité corrigé, exemple parfait de traçabilité de conception. Signaler aussi
l'absence de \fn{require\_role} sur ce routeur.

% =====================================================================
\section{\file{backend/app/services/incident\_service.py}}

\textbf{Objectif} --- Toutes les écritures liées aux incidents et aux actions.

\textbf{Fonctions} --- \fn{create\_action\_request},
\fn{confirm\_and\_execute\_action}, \fn{reject\_action\_request},
\fn{\_generate\_ticket\_number}, \fn{\_open\_ticket\_for},
\fn{create\_ticket\_from\_diagnosis}, \fn{list\_tickets}, \fn{get\_ticket},
\fn{update\_ticket\_status}, \fn{mark\_in\_progress}, \fn{add\_note},
\fn{resolve\_ticket}, \fn{pending\_actions}, \fn{has\_completed\_action}.

\textbf{Fichiers qui l'appellent} --- \file{agent.py}, \file{actions.py},
\file{incidents.py}, \file{incident\_investigation.py},
\file{incident\_watcher.py}, \file{report\_service.py}.

\textbf{À retenir} --- Trois protections en dix lignes dans
\fn{confirm\_and\_execute\_action} : anti-double-exécution, commit de
l'intention avant l'action, et détection d'une erreur retournée plutôt que
levée.

% =====================================================================
\section{\file{backend/app/services/incident\_watcher.py}}

\textbf{Objectif} --- Détecter les échecs de sauvegarde sans intervention
humaine, et refermer les tickets vérifiés.

\textbf{Fonctions} --- \fn{\_system\_user\_id}, \fn{\_parse\_timestamp},
\fn{\_latest\_ticket}, \fn{\_already\_handled}, \fn{\_verify\_resolved},
\fn{scan\_once}, \fn{run\_forever}.

\textbf{Constantes} --- \texttt{INTERVAL\_SECONDS = 300},
\texttt{SYSTEM\_ACCOUNT\_EMAIL = "oasis-system@oasis.local"}.

\textbf{Sorties} --- Un rapport
\texttt{\{"created", "skipped", "closed", "checked"\}} ; jamais d'exception.

\textbf{À retenir} --- Le compte de service désactivé est une élégance :
l'audit distingue « ce qu'OASIS a détecté seul » de « ce qu'un humain a
demandé », sans créer de faille de connexion.

% =====================================================================
\section{\file{backend/app/dashboard/kpi\_service.py}}

\textbf{Objectif} --- Produire, sans LLM, tout ce qu'affiche le dashboard.

\textbf{Fonctions} --- \fn{\_volume\_status}, \fn{\_build\_capacity},
\fn{\_build\_backups}, \fn{\_build\_incidents}, \fn{\_health\_score},
\fn{\_headline}, \fn{refresh\_health\_snapshot}, \fn{\_storage\_health},
\fn{\_refresh\_in\_background}, \fn{\_invalidate\_health\_snapshot},
\fn{get\_overview}, \fn{get\_performance}, \fn{get\_recent\_activity},
\fn{get\_insights}.

\textbf{État partagé} --- \texttt{\_health\_cache} protégé par deux verrous
(\texttt{\_health\_lock} pour la lecture/écriture, \texttt{\_refresh\_lock} pour
sérialiser les rafraîchissements).

\textbf{À retenir} --- Le cache de 60 secondes avec rafraîchissement en arrière-
plan est la réponse à un problème mesuré : deux connexions SSH à chaque
affichage de page.

% =====================================================================
\section{\file{backend/app/rag/embedding\_service.py} et \file{vectorstore.py}}

\textbf{Objectif} --- Vectoriser, indexer et rechercher.

\textbf{Fonctions} --- \fn{get\_embedding\_model} (chargement paresseux, une
fois par processus), \fn{embed\_query}, \fn{embed\_passages} ;
\fn{get\_chroma\_client}, \fn{get\_or\_create\_collection}, \fn{add\_chunks},
\fn{search}.

\textbf{Classes} --- \cls{EmbeddingModelMismatch(RuntimeError)}.

\textbf{À retenir} --- Deux garde-fous rarement vus dans un projet d'école : la
séparation stricte query/passage imposée par l'asymétrie des modèles E5, et le
refus d'une collection indexée avec un autre modèle --- parce que deux modèles
de même dimension produiraient « des réponses fausses, mais parfaitement
plausibles ».

% =====================================================================
\section{\file{frontend/src/lib/api.ts}}

\textbf{Objectif} --- Client HTTP unique du frontend.

\textbf{Objets exportés} --- \cls{ApiError}, \fn{tokenStore},
\fn{setUnauthorizedHandler}, \fn{api} (\fn{get}, \fn{post}, \fn{patch},
\fn{delete}, \fn{upload}, \fn{download}).

\textbf{Fichiers qui l'appellent} --- \cls{AuthContext} et les huit pages.

\textbf{À retenir} --- \fn{api.download} explique bien un problème réel : « un
simple \texttt{<a href>} ne convient pas : le navigateur n'y joint pas
l'en-tête \texttt{Authorization} ». D'où le passage par un blob et une URL
objet.

% =====================================================================
\section{\file{frontend/src/components/ActionProposal.tsx}}

\textbf{Objectif} --- Matérialiser la confirmation humaine.

\textbf{Fonctions} --- \fn{describe}, \fn{resultText}, \fn{settle},
composant \cls{ActionProposal}.

\textbf{Entrées} --- \texttt{action: ActionRequest},
\texttt{onSettled?: (updated) => void}.
\textbf{Sorties} --- Appels \route{POST /storage/actions/\{id\}/confirm} ou
\route{/reject}, notifications \emph{toast}.

\textbf{Fichiers qui l'appellent} --- \cls{ChatPage} et \cls{IncidentsPage}.

\textbf{À retenir} --- Le composant est \emph{partagé} entre le chat et la fiche
d'incident, « pour que la décision se prenne au même endroit que la lecture,
quel que soit l'écran ». Il affiche à la fois la phrase lisible et l'appel
technique brut \texttt{action\_type(target)}, ce qui permet de vérifier la cible
avant de cliquer.

% =====================================================================
\section{\file{frontend/src/pages/IncidentsPage.tsx}}

\textbf{Objectif} --- Liste, fiche détaillée et pilotage du cycle de vie.

\textbf{Constantes} --- \texttt{FILTERS} (all, critical, high, medium, low,
resolved), \texttt{LIFECYCLE} (5 étapes).

\textbf{Fonctions} --- \fn{queryFor}, \fn{reachedStage}, \fn{Lifecycle},
\fn{IncidentBody}, \fn{DetailRow}.

\textbf{À retenir} --- \texttt{LIFECYCLE} n'est pas décoratif. Le commentaire du
fichier le dit : « L'ordre n'est pas décoratif : c'est le contrat de
l'application. Rien ne s'exécute sur la VM sans être passé par "Action proposée"
puis par le clic humain de "Confirmation", et un ticket ne se referme qu'une
fois l'effet vérifié. »

\begin{lstlisting}[caption={Les cinq temps affichés, et la fonction qui calcule l'étape atteinte}]
const LIFECYCLE = [
  { key: 'detected',  label: 'Detected' },
  { key: 'diagnosed', label: 'Diagnosed' },
  { key: 'proposed',  label: 'Action proposed' },
  { key: 'confirmed', label: 'Confirmed' },
  { key: 'resolved',  label: 'Resolved' },
];

function reachedStage(incident) {
  if (incident.status === 'resolved' || incident.status === 'closed') return 4;
  if (incident.actions.some((a) => a.status === 'completed'))          return 3;
  if (incident.actions.length > 0)                                     return 2;
  if (incident.root_cause)                                             return 1;
  return 0;
}
\end{lstlisting}

% =====================================================================
\chapter{Annexes}

\section{Fichiers présents mais vides}

\begin{longtable}{@{}p{7.4cm}p{7.4cm}@{}}
\toprule
\textbf{Fichier} & \textbf{Ce que son nom laisse attendre}\\
\midrule
\endhead
\file{backend/Dockerfile} & Image du backend\\
\file{frontend/Dockerfile} & Image du frontend\\
\file{nginx/nginx.conf} & Reverse proxy\\
\file{.github/workflows/ci.yml} & Intégration continue\\
\file{backend/app/core/permissions.py} & Matrice de permissions\\
\file{backend/app/core/constants.py} & Constantes partagées\\
\file{backend/app/database/init\_db.py} & Initialisation du schéma\\
\file{backend/app/orchestrator/orchestrator.py} & Orchestrateur multi-agents\\
\file{backend/app/orchestrator/cross\_domain\_manager.py} & Gestion cross-domain\\
\file{backend/app/llm/models.py}, \file{prompt\_manager.py} & Abstractions LLM\\
\file{backend/app/middleware/auth\_middleware.py} & Middleware d'authentification\\
\file{backend/app/reports/generator.py} & Générateur de rapports\\
\file{backend/app/models/report.py}, \file{llm\_usage.py} & Modèles non créés\\
\file{backend/app/schemas/} (plusieurs) & \file{audit\_log.py}, \file{llm\_usage.py}, \file{mcp\_call.py}, \file{metric.py}, \file{report.py}, \file{role.py}\\
\bottomrule
\end{longtable}

\begin{infobox}
Ces fichiers ne gênent rien : ils ne sont importés par aucun module actif. Ils
témoignent d'une architecture cible plus vaste que ce qui a été réalisé --- ce
qui est normal pour un projet de stage, et vaut mieux dit qu'esquivé.
\end{infobox}

\section{Code présent mais non appelé}

\begin{longtable}{@{}p{5.4cm}p{9.4cm}@{}}
\toprule
\textbf{Élément} & \textbf{Observation}\\
\midrule
\endhead
\file{personas/base\_persona.py} & \cls{BasePersona} n'est héritée par personne ; \cls{StoragePersona} ne l'étend pas\\
\file{orchestrator/llm\_router.py} & \fn{route()} lève \texttt{NotImplementedError} pour \texttt{COMPLEX} ; appelée uniquement par \cls{BasePersona}\\
\file{middleware/rbac\_middleware.py} & \fn{require\_role} en double de celle de \file{api/deps.py} ; c'est la seconde qui est importée partout\\
\fn{risk\_for}, \fn{invoke} & Définies dans \file{mcp/storage/server.py}, jamais appelées\\
\fn{get\_calls\_for\_correlation} & Définie dans \file{mcp\_audit.py}, jamais appelée\\
\cls{IncidentAnalysisRequest} / \cls{IncidentAnalysisResponse} & Schémas définis, aucune route ne les utilise\\
\tbl{request\_messages} & Table conservée pour l'historique, plus alimentée\\
\bottomrule
\end{longtable}

\section{Comptes de test}

Créés par \texttt{python -m app.database.seed} :

\begin{longtable}{@{}p{5.0cm}p{4.0cm}p{5.2cm}@{}}
\toprule
\textbf{E-mail} & \textbf{Mot de passe} & \textbf{Rôle}\\
\midrule
\endhead
\texttt{linux@oasis.com} & \texttt{linux123} & \texttt{linux\_engineer}\\
\texttt{storage@oasis.com} & \texttt{storage123} & \texttt{storage\_engineer}\\
\texttt{admin@oasis.com} & \texttt{admin123} & \texttt{administrator}\\
\bottomrule
\end{longtable}

\begin{warnbox}
Ces identifiants sont en clair dans \file{app/database/seed.py}. Acceptable pour
une maquette, à retirer avant tout déploiement.
\end{warnbox}

\section{Démarrage de l'environnement}

\begin{lstlisting}[caption={Séquence complète, d'après le code et .claude/launch.json}]
# 1. Dependances (racine du depot)
docker compose up -d                 # postgres:16 + chromadb

# 2. Backend
cd backend
alembic upgrade head                 # 14 migrations
python -m app.database.seed          # roles + comptes de test
python -m app.rag.reindex            # (re)construction des collections Chroma
uvicorn app.main:app --reload --port 8000

# 3. Frontend
cd frontend
npm install
npm run dev                          # http://localhost:5173, proxy /api -> :8000
\end{lstlisting}

\section{Dix points à retenir pour la soutenance}

\begin{enumerate}
  \item \textbf{Rien n'est simulé.} Chaque chiffre affiché vient d'une commande
  Linux exécutée par SSH sur une vraie VM.
  \item \textbf{Le modèle propose, l'humain décide.} Quatre outils sur dix ne
  s'exécutent jamais depuis le chat.
  \item \textbf{Les seuils sont dans le code, pas dans le modèle.}
  \file{incident\_rules.py} et \file{health.py} décident ce qui est un incident.
  \item \textbf{Un ticket ne se ferme que sur une vérification.}
  \file{incident\_verification.py} relit l'indicateur qui l'avait ouvert.
  \item \textbf{Tout est audité.} \tbl{mcp\_calls}, \tbl{audit\_logs},
  \tbl{incident\_ticket\_notes}, \tbl{request\_events}.
  \item \textbf{Le RAG est multilingue par nécessité.} KB en anglais, questions
  en français ; d'où E5 multilingue et ses préfixes asymétriques.
  \item \textbf{MCP isole les outils dans un processus séparé} et fournit le
  catalogue que le modèle lit pour choisir.
  \item \textbf{Trois filets empêchent de perdre une remédiation} : appel
  d'outil, rattrapage d'un appel écrit en prose, action de repli calculée par la
  règle.
  \item \textbf{Le système travaille seul.} Deux boucles de fond détectent,
  mesurent et referment sans qu'on le leur demande.
  \item \textbf{Les limites sont connues et nommées} : CORS ouvert, absence de
  \fn{require\_role} sur les actions, jeton cross-domain non vérifié,
  conteneurisation incomplète, Claude non branché.
\end{enumerate}

\begin{keybox}
Le dernier point est le plus fort. Un projet dont l'auteur sait exactement où
sont les limites inspire davantage confiance qu'un projet présenté comme
parfait.
\end{keybox}
